\documentclass[11pt]{article}
\usepackage{amsmath,amssymb,amsthm}
\usepackage[margin=1in]{geometry}

\title{Theorem 5: A Density-Regularity Clarification}
\author{}
\date{}

\begin{document}
\maketitle

\section*{Statement and sufficient replacement}

Theorem 5 of Kleinberg and Raghavan's \emph{Algorithmic Monoculture and Social
Welfare} (arXiv:2101.05853, version 2, Appendix A)
claims that a differentiable density $f$ with positive support on all of
$\mathbb R$ makes every scaled-noise ranking atom differentiable in the
positive accuracy parameter $\theta$.  The printed proof differentiates an
integrand pointwise and moves that derivative through an unbounded integral.

A smooth positive density can have increasingly narrow peaks that prevent
ranking probabilities from being differentiable. The construction below has
finite variance and can be normalized to mean zero and variance one. A
sufficient replacement is global absolute continuity of the density with
an integrable derivative, $f\in W^{1,1}(\mathbb R)$. This condition covers
Gaussian and Laplace noise; its necessity is not asserted.

\section*{Construction}

Fix $a=1$, and choose a nonnegative nonzero
$\psi\in C_c^\infty((-1,1))$.  For $k\geq 2$, let
\[
  c_k=4k,\qquad h_k=k^6,\qquad w_k=\frac{1}{100k^{12}},
\]
and define the paired bumps
\[
  b_k(x)=h_k\psi\left(\frac{x-c_k}{w_k}\right),\qquad
  b'_k(x)=h_k\psi\left(\frac{x-(c_k+a)}{w_k}\right).
\]
Their supports are pairwise disjoint. Put
\[
  q(x)=e^{-x^2}+\sum_{k\geq2}(b_k(x)+b'_k(x)),
  \qquad f(x)=\frac{q(x)}{\int_{\mathbb R}q(u)\,du}.
\]

The sum is locally finite, hence $q$ and $f$ are $C^\infty$.  The Gaussian
summand makes $f(x)>0$ for every real $x$.  Moreover,
\[
  \int b_k = h_kw_k\int\psi=O(k^{-6}),
  \qquad
  \int x^2b_k(x)\,dx=O(k^2h_kw_k)=O(k^{-4}),
\]
and the same bounds hold for $b'_k$.  Thus $f$ is a positive smooth probability
density with finite second moment.  Centering and rescaling it to unit variance
preserves all of these properties and only rescales the positive gap $a$.

\section*{Failure of ranking-atom differentiability}

Let $\epsilon_1,\epsilon_2$ be iid with density $f$, and let
$Z=\epsilon_2-\epsilon_1$.  Its density is
\[
  g(t)=\int_{\mathbb R}f(x)f(x+t)\,dx.
\]
At $t=a$, the nonnegative paired-bump terms give
\[
  g(a)\ \geq\ C\sum_{k\geq2}
    h_k^2w_k\int\psi(u)^2\,du
  = C'\sum_{k\geq2}1=+\infty,
\]
where $C,C'>0$ include the normalization factor.  More strongly, if
$\kappa(r)=\int\psi(u)\psi(u+r)\,du$, then $\kappa$ is continuous and
$\kappa(0)>0$.  Hence there are $\eta,\gamma>0$ such that
$\kappa(r)\geq\gamma$ for $|r|\leq\eta$.  For $0<r$ sufficiently small,
every $k$ with $w_k\geq r/\eta$ contributes at least a fixed positive
constant throughout $[a,a+r]$.  Therefore
\[
  \frac{\Pr[a<Z\leq a+r]}{r}
  =\frac1r\int_a^{a+r}g(t)\,dt
  \ \longrightarrow\ +\infty.
\]
The distribution function of $Z$ consequently has no finite right derivative
at $a$.

Now take two candidate values $x_1>x_2$ and a positive $\theta_0$ with
$\theta_0(x_1-x_2)=a$.  Under the source RUM convention, the atom that ranks
candidate $1$ first is
\[
  p(\theta)=
    \Pr\left[x_1+\frac{\epsilon_1}{\theta}>
               x_2+\frac{\epsilon_2}{\theta}\right]
    =\Pr[Z<\theta(x_1-x_2)].
\]
Its right difference quotient at $\theta_0$ is the divergent quotient above,
multiplied by the positive constant $x_1-x_2$.  Thus $p$ is not differentiable
at a positive accuracy, although $f$ satisfies the printed density hypotheses.

\section*{Why the sufficient condition works}

The formalized statement uses $f\in W^{1,1}(\mathbb R)$ together with the
source's normalization and positive-support conditions. For $\theta>0$,
\[
  \operatorname{rank}\left(x_i+\frac{\epsilon_i}{\theta}\right)
  =\operatorname{rank}(\theta x_i+\epsilon_i).
\]
For a fixed ranking cell $C_\pi$, this gives
\[
  p_\pi(\theta)=\int_{C_\pi}\prod_i f(r_i-\theta x_i)\,dr.
\]
Translation is differentiable in $L^1$ under $W^{1,1}$, with derivative
$-x_i f'(r_i-\theta x_i)$.  A finite telescoping product argument gives an
$L^1$ derivative of the joint density, and integration over $C_\pi$ is a
bounded linear functional on $L^1$.  This yields differentiability of every
ranking atom without a pointwise common envelope.

The $W^{1,1}$ ranking result is proved in Lean. The counterexample above is a
mathematical argument supplied in this note, rather than a Lean-checked theorem.

\end{document}
